% paper/drafts/open_problem.tex -- DRAFT, standalone section, meant for \input from main.tex.
% Status: pending. Every statement below repeats FABLE_HARD.md, which has no verifier report yet
% (FABLE_HARD_VERIFY.md does not exist on 2026-10-28). Do not cite as verified.
%
% Paper round 10 (2026-10-28): the second problem is now the K-uniform (Gap_U) where it concerns OS4
%   for joint limits, and Theorem U1 is replaced by the one-way form U1' ((U-Lap)_K <=> (U-IV=)_K;
%   (U-Lap)_K => (U-IV)_K one-way), both per SOURCES_FH_U1_REPAIR.md \S2.3, \S3.2-\S3.5
%   (an Opus repair file, UNVERIFIED). Still pending: the FABLE_HARD verifier (OPUS_QUEUE VH1).
% IDS: the ledger row B18 (N3) is now written R-042 (N3) (IDS.md). File-local labels are kept and
%   qualified by file: L3 (FX-V) = L-320 / H-106, L6 (S1-pt) = L-323 / H-165 with forms (L6-nu), (L6-P),
%   L14 (Loc-f) = L-331 / H-142, Theorem IV = L-196 (Theorem~\ref{thm:L196} of main.tex), E13, (Q),
%   Theorems U1, U1', U2, W1, W1_fin, Corollaries W2, W2_J, Theorem FD, Proposition RP1 (no IDS rows).
%   The labels sec:open, prob:ULap, prob:Gap, ... are provisional and follow no IDS scheme yet.
%
% Sources used (file : lines), read 2026-10-28:
%   FABLE_HARD.md : 1-17 (orientation, grades legend), 28-34 (L6-nu, L6-P, polydisc P),
%                   113-138 (Task 3: objects, Theorem U1, Theorem U2, joint limit, verdict),
%                   150-151 (Theorem FD, failing inequality), 155-158 (Proposition RP1, ergodic decomposition),
%                   160-170 (Theorem W1, Corollary W2, (Gap_K), verdict), 172-184 (summary table).
%   INFINITE_VOLUME_FABLE.md : 54-55 (Step B not l-uniform), 57-61 (nu_{K,l}, Delta, (2.1)),
%                   102-124 (Theorem LOC, (3.1), (3.2), compactness), 186-216 (Theorem IV, lemmas L1-L17).
%   PLAN.md : 551 (FABLE_HARD status: UNVERIFIED, verifier launched), 552 (ALT_ROUTES restatement of (Gap_K)).
%   SOURCES_FH_U1_REPAIR.md : 9-16 (results), 76-85 (Theorem U1'), 86-131 (W1_fin, W2_J, (Gap_U), \S3.6),
%                   145-175 (L6 grade, mapping to this draft). Added in paper round 10.
%   paper/main.tex : 1-36 (preamble: environments, \status, \file, \ioc).
%   paper/remains.tex : 91-93 (Section rem-gap, the regime separation).
%   IDS.md : 1-12, 369, 405, 428, 802 (id scheme, sample mappings).
% Firewall: /tmp/sa_repo, reports/SUSSMAN_ALPOGE_MAP.md, reports/ONLINE_SWEEP.md not read.
%
% The preamble of main.tex has no environment for open problems. This guard defines one;
% move it to the preamble when the draft is merged.
\ifdefined\problem\else
\theoremstyle{definition}
\newtheorem{problem}[theorem]{Open problem}
\fi

\section{Two stated open problems}\label{sec:open}

\noindent\textbf{[Pending verification: depends on FABLE\_HARD verifier and on a verifier of SOURCES\_FH\_U1\_REPAIR]}

\medskip
This section states two hypotheses as open problems. Neither is proved in this paper or in the project files. Both concern the infinite-volume limit $\ell\to\infty$, which the main theorems of this paper do not treat (Section~\ref{sec:notgiven}). The first, (U-Lap)$_K$, is a sufficient condition for uniqueness of the infinite-volume limit at fixed lattice spacing; the converse is not proved. The second, (Gap$_U$), is a sufficient condition for clustering (OS4) of the joint limits $(K,\ell)\to(\infty,\infty)$ of Theorem IV and, in its exponential form, for a mass gap. It is uniform in $K$. The fixed-$K$ condition (Gap$_K$) gives clustering only for the infinite-volume lattice limits at one fixed $K$. The source of both problems is \file{FABLE_HARD.md} \S3--\S4, as repaired by \file{SOURCES_FH_U1_REPAIR.md} \S2--\S3, built on the infinite-volume theorem of \file{INFINITE_VOLUME_FABLE.md} \S7.1 (Theorem IV, L-196). All results quoted here carry the grades of those two files; no verifier has reported on either.

\subsection{Setting}\label{sec:open-setting}

Fix $G=\mathrm{SU}(2)$, $d=4$, odd $L>11$, and $g\in\ioc{\gamma}$ with $\gamma$ as in L-001 (Theorem 2w). Fix the number $K$ of renormalization-group steps, so the lattice spacing $\varepsilon=L^{-K}$ is fixed. Let $T_\ell$ be the unit-scale torus of even side $\ell>22M$.

\begin{definition}[unit-scale measure and response functional; \file{INFINITE_VOLUME_FABLE.md} \S2.2, \S3.1]\label{def:open-nu}
Let $\hat\rho^f_K$ be the unit-scale density with source $f$ after Step~A of (FX). The measure
\[
\nu_\ell:=\nu_{K,\ell}(dV):=\hat\rho^0_K(V)\,dV\big/I_K[0]
\]
is a probability measure on $\mathrm{SU}(2)^{\{\text{unit links of }T_\ell\}}$; positivity of $\hat\rho^0_K$ is Ba{\l}aban's positivity of $\rho_K$. Put $\Delta(f,V):=\log\hat\rho^f_K(V)-\log\hat\rho^0_K(V)$. Then $Z_\varepsilon[f]/Z_\varepsilon[0]=e^{-W_K[f]}\,\mathbb E_{\nu_\ell}[e^{\Delta(f,\cdot)}]$ exactly.
\end{definition}

Let $f_1,\dots,f_n$ be bumps with disjoint supports. Write $N(f_i):=\|f_i\|_\infty+\|\nabla f_i\|_\infty$ and $|\operatorname{supp}f_i|_1$ for the number of unit cubes that meet $\operatorname{supp}f_i$. The common source polydisc is
\[
\mathcal P:=\Bigl\{t\in\mathbb C^n:\ |t_i|\le c_*r/N(f_i),\ i=1,\dots,n\Bigr\},\qquad f(t):=\textstyle\sum_it_if_i,
\]
with $c_*r$ the common radius of \file{SOURCES_FX_D.md} \S5.2; this radius needs R-042 (N3). Put $\Delta_\ell(t,V):=\Delta(f(t),V)$ on $T_\ell$. Let $\Delta_\infty(t,V)$ be the same sum of polymer activities built from the $\mathbb Z^4$ activities, with non-wrapping polymers only. The centring vector is $m_\ell:=\bigl(\mathbb E_{\nu_\ell}[\partial_{t_i}\Delta_\ell(0,\cdot)]\bigr)_{i\le n}$. The response family is
\[
\mathfrak F:=\bigl\{\Delta_\infty(t,\cdot):\ f_1,\dots,f_n\text{ a disjoint bump family},\ t\in\mathcal P\bigr\},
\]
and the first-order response of a bump $f_i$ is $w_i:=\partial_{t_i}\Delta_\infty(0,\cdot)$.

Three inputs of Theorem IV are used throughout: L3 (FX-V), the polymer representation of $\Delta$ uniform in $V$, $K$, $\ell$; L6 in its polydisc form (L6-$\mathcal P$),
\[
\sup_V\bigl|\Delta(f(t),V)-\Delta(f(t),1)\bigr|\le C_2\sum_i|t_i|\,N(f_i)\,(1+|\operatorname{supp}f_i|_1)/r\qquad(t\in\mathcal P),
\]
with $C_2$ free of $K$, $\ell$, $V$; and L14 (Loc-f), which gives $\sup_{t,V}|\Delta_\ell-\Delta_\infty|\le Ce^{-c\ell}$ (\file{INFINITE_VOLUME_FABLE.md} (3.1)). \file{FABLE_HARD.md} \S1 grades (L6-$\mathcal P$) \status{proved relative to I1--I8 and R-042, unverified}; \file{SOURCES_FH_U1_REPAIR.md} \S4.2 regrades it \status{P-sketch} relative to I1--I8, R-042, and vacuum extraction (unverified). L3 and L14 are \status{reductions, unverified} (\file{INFINITE_VOLUME_FABLE.md} \S7.2).

\subsection{Uniqueness at fixed lattice spacing: (U-Lap)$_K$}\label{sec:open-ulap}

\begin{problem}[(U-Lap)$_K$; \file{FABLE_HARD.md} \S3.2]\label{prob:ULap}
Fix $K$. Prove that for every $n\ge1$, every family $f_1,\dots,f_n$ of bumps with disjoint supports, and every $t\in\mathcal P$, the limit
\[
\lim_{\ell\to\infty}\ \mathbb E_{\nu_\ell}\Bigl[\exp\bigl(\Delta_\infty(t,\cdot)-t\cdot m_\ell\bigr)\Bigr]
\]
exists, with $\ell$ running through the even integers $\ell>22M$.
\end{problem}

The functionals $\Delta_\infty(t,\cdot)$ do not depend on $\ell$. Only the measure $\nu_\ell$ and the centring vector $m_\ell$ depend on $\ell$. So (U-Lap)$_K$ is a statement about the unit-scale measures alone: convergence of their centred exponential moments on $\mathfrak F$.

\paragraph{Where it enters.} Theorem IV (\file{INFINITE_VOLUME_FABLE.md} \S7.1) gives a joint subsequential limit $(K_i,\ell_i)\to(\infty,\infty)$ of the Schwinger functions of the centred plaquette density. Its part (d) does not assert convergence along the full sequence. Lemma L15 (full convergence in $\ell$) is graded \status{blocked: infrared} there. (U-Lap)$_K$ is a sufficient replacement for L15 at fixed $K$. It is not an equivalent one. The statement below is the corrected form of Theorem U1 of \file{FABLE_HARD.md} \S3.2. Write (U-IV$^=$)$_K$ and (U-cum$^=$)$_K$ for the forms of (U-IV)$_K$ and (U-cum)$_K$ in which coincident points (repeated plaquettes) are allowed.

\begin{theorem}[Theorem U1$'$ of \file{SOURCES_FH_U1_REPAIR.md} \S2.3 \status{proved given (L6-$\mathcal P$), L3, L14; unverified}]\label{thm:open-U1}
Fix $K$. Let $\mathcal P^\circ$ be the open polydisc $\mathcal P$; its radius needs R-042. Consider the statements
\begin{enumerate}
\item[(U-IV)$_K$] $\lim_{\ell\to\infty}S_{n,\varepsilon,\ell}(F)$ exists for every $n$ and every $F\in\bigcup_BD_n(B)$;
\item[(U-cum)$_K$] $\lim_{\ell\to\infty}\kappa_B^{(K,\ell)}(f_1,\dots,f_{|B|})$ exists for every $B$ and every disjoint bump family.
\end{enumerate}
Then:
\begin{enumerate}
\item[(a)] (U-Lap)$_K$ on $\mathcal P^\circ$ $\Leftrightarrow$ (U-cum$^=$)$_K$ $\Leftrightarrow$ (U-IV$^=$)$_K$;
\item[(b)] (U-IV)$_K$ $\Leftrightarrow$ (U-cum)$_K$;
\item[(c)] hence (U-Lap)$_K$ $\Rightarrow$ (U-IV)$_K$. This implication is one-way. The converse holds if and only if a statement (U-diag)$_K$ on coincident points holds in addition; (U-diag)$_K$ is not proved, and the converse fails for general families of measures.
\end{enumerate}
\end{theorem}

The proof in the repair file has three steps. First, centred moments and cumulants of order $\ge2$ determine each other by finite partition sums, and bump families are dense in $D_n$; this gives (b). Second, put $\Phi_\ell(t):=\mathbb E_{\nu_\ell}[\exp(\Delta_\ell(t,\cdot)-t\cdot m_\ell)]$. By (L6-$\mathcal P$), $(\Phi_\ell)$ is a uniformly bounded family of holomorphic functions on $\mathcal P^\circ$, and its Taylor coefficients at $0$ are polynomials in the cumulants, with repeated indices included. Vitali's theorem, applied to $\Phi_\ell$ and not to $\log\Phi_\ell$, gives (a), after the replacement $\Delta_\ell\to\Delta_\infty$ at cost $e^{Ce^{-c\ell}}-1$. No lower bound on $|\mathbb E\,e^{\Delta^c}|$ is used. (\file{FABLE_HARD.md} l.~118 applies Vitali to $\log\Phi_\ell$ and asserts the range $[e^{-C_{\mathcal P}},e^{C_{\mathcal P}}]$. A sup-based lower bound needs $C_1<\pi/2$, and it fails on $\mathcal P$ for $n>\pi/(4c_*C_2)$; \file{SOURCES_FH_U1_REPAIR.md} \S1.3--\S1.4.) Third, the coefficients with a repeated index are cumulants at coincident plaquettes. Test functions in $D_n$ vanish on the diagonals, so (U-IV)$_K$ does not control these coefficients. This gives (c).

So, by (c), if (U-Lap)$_K$ holds, the limit $\ell\to\infty$ of Theorem IV exists at that fixed $K$ along the full sequence of even $\ell$, not only along a subsequence. This is the ``Theorem IV-unique'' variant that \file{FABLE_HARD.md} \S3.5 proposes to record.

\paragraph{What it does not give.} (U-Lap)$_K$ for every $K$ does not by itself make the joint limit $(K,\ell)\to(\infty,\infty)$ unique (\file{FABLE_HARD.md} \S3.4). Two further inputs are needed: convergence of the fixed-$K$ limits $S^{(K,\infty)}$ as $K\to\infty$ (an E13-type ultraviolet uniqueness statement, open), and a uniformity statement, for example that the $\ell$-convergence in (U-Lap)$_K$ is uniform in $K$ on each bump family. Iterated limits do not give a joint limit without it. The precise form is this: if (U-IV)$_K$ holds for every $K\ge K_0$ with $S_{K,\ell}(F)\to S^{(K)}(F)$ uniformly in $K\ge K_0$, and $S^{(K)}(F)\to S(F)$ as $K\to\infty$, then $S_{K,\ell}(F)\to S(F)$ along every joint sequence (\file{SOURCES_FH_U1_REPAIR.md} \S3.6, unverified). Only (U-IV)$_K$ is needed here. The converse claim of \file{FABLE_HARD.md} \S3.4, that uniqueness of the joint limit along all sequences implies (U-Lap)$_K$ for every $K$, is false: $S_{K,\ell}=S+K^{-1}(-1)^{\ell/2}$ has the limit $S$ along every joint sequence and no $\ell$-limit at any fixed $K$. The chain of this paper controls none of the three inputs.

\paragraph{What is known.} \file{FABLE_HARD.md} \S3.5 grades the item \status{equivalent (fixed $K$)}; \file{SOURCES_FH_U1_REPAIR.md} \S2.3 changes this to \status{equivalent to (U-IV$^=$)$_K$; sufficient for (U-IV)$_K$}. Neither file gives a proof route. The known facts are these.
\begin{enumerate}
\item The $\ell$-dependence of every Schwinger function, beyond $O(e^{-c\ell})$, sits in the expectations and joint cumulants under $\nu_\ell$ of fixed local functionals (\file{INFINITE_VOLUME_FABLE.md} \S3.2). The ultraviolet activities are fixed functions; only the measure moves.
\item Every subsequential weak limit $\nu_\infty$ of $(\nu_\ell)$ is a translation-invariant, $B_4$-invariant Gibbs measure for an absolutely summable interaction $\Phi_K$ on $\mathbb Z^4$ (Theorem U2(a) of \file{FABLE_HARD.md} \S3.3, \status{proved given L14 and (Q), unverified}). Here $\Phi_K$ consists of the plaquette terms $-g^{-2}A_p$, the single-plaquette factors of $\chi_K$, the unit-scale terms $E^{(K)}_{\mathbb Z^4}$ and $R_{\mathrm{fin},\mathbb Z^4}$, and the transported fine terms $E^{(j)}$, $R^{(j)}$ evaluated at $U_j(\cdot)$. Hypothesis (Q) is quasi-locality of the transported terms (\status{R}: A1-type kernel decay of the minimizer).
\item The interaction $\Phi_K$ is weakly coupled ($g\le\gamma$). Its expected correlation length is $\xi_g\sim e^{c/g^2}$. Uniqueness of its Gibbs state is the infrared problem of four-dimensional lattice gauge theory at weak coupling. No tool of the project or of Ba{\l}aban's papers addresses it. Griffiths and FKG inequalities are not available for SU(2) (\file{FABLE_HARD.md} \S4.3).
\item On a finite window of volumes, the Step-B estimates hold with $\ell$-free constants (Theorem FD of \file{FABLE_HARD.md} \S4.2, \status{proved given (c$'$) for the continued flow, unverified}): for every even $\ell\le\ell_*(g):=22ML\,\xi_g$, with $\xi_g=L^{k_*(g)}=\exp((g^{-2}-\gamma^{-2})/c+O(1))$, $c\approx0.3715$ (project value). This window is finite for each $g$. It gives no statement about the limit $\ell\to\infty$.
\end{enumerate}

\begin{remark}[equivalent, stronger, and weaker forms of (U-Lap)$_K$]\label{rem:open-ulap-forms}
\begin{enumerate}
\item Equivalent forms. By Theorem~\ref{thm:open-U1}, (U-Lap)$_K$ is equivalent to (U-cum$^=$)$_K$ and to (U-IV$^=$)$_K$ at fixed $K$, and it implies (U-IV)$_K$. \file{FABLE_HARD.md} \S3.3(c) calls (U-Lap)$_K$ the exact equivalence class of (U-IV)$_K$; the repair file corrects this to (U-IV$^=$)$_K$.
\item A stronger, sufficient form. Let (U-Gibbs)$_K$ be the statement ``$\Phi_K$ has a unique translation-invariant Gibbs measure''. Then (U-Gibbs)$_K$ implies $\nu_\ell\to\nu_\infty$ weakly, which implies (U-Lap)$_K$, which implies (U-IV)$_K$ (Theorem U2(b), \status{proved given L14 and (Q), unverified}).
\item The converse is not claimed. The family $\mathfrak F$ need not separate Gibbs measures of $\Phi_K$. Its first-order part is the smeared unit-scale plaquette action $B_K\sum_pf(x_p)A_p(V)$ plus $E$-corrections, and plaquette traces do not generate the full gauge-invariant local algebra (a $1\times2$ loop is not a function of plaquette traces). So (U-Lap)$_K$ lies strictly between ``no condition'' and (U-Gibbs)$_K$ as stated in the source; whether it is strictly weaker than (U-Gibbs)$_K$ is not settled.
\item The sources give no weaker form that still implies (U-IV)$_K$.
\end{enumerate}
\end{remark}

\subsection{Clustering and the mass gap: (Gap$_U$)}\label{sec:open-gap}

For bumps $f_1,f_2$ and $a\in\mathbb R^4$, write $w_1:=\partial_{t_1}\Delta_\ell(0,\cdot)$ for $f_1$ and $w_{2,a}:=\partial_{t_2}\Delta_\ell(0,\cdot)$ for the translate $\tau_af_2$, on $T_\ell$ at level $K$. These are the finite-volume first-order responses; the $w_i$ of Section~\ref{sec:open-setting} are their $\mathbb Z^4$ forms.

\begin{problem}[(Gap$_U$); \file{SOURCES_FH_U1_REPAIR.md} \S3.5]\label{prob:Gap}
Prove that for every pair of disjoint bumps $f_1,f_2$,
\[
\lim_{R\to\infty}\ \limsup_{K,\ell\to\infty}\ \sup_{R\le|a|\le\ell/4}\bigl|\operatorname{Cov}_{\nu_{K,\ell}}(w_1,w_{2,a})\bigr|=0 .
\]
The exponential form, (Gap$_U$) with rate $m$, asks for $|\operatorname{Cov}_{\nu_{K,\ell}}(w_1,w_{2,a})|\le A\,e^{-m|a|}$ for all $K\ge K_0$ and all $\ell\ge4|a|$, with $A=A(f_1,f_2)$ and some $m>0$.
\end{problem}

(Gap$_U$) is a statement about the finite-volume unit-scale measures $\nu_{K,\ell}$, uniform in $K$ and $\ell$. Where OS4 for the joint limits of Theorem IV is concerned, it replaces the fixed-$K$ statement of \file{FABLE_HARD.md} \S4.4:

\begin{definition}[(Gap$_K$); \file{FABLE_HARD.md} \S4.4]\label{def:open-GapK}
Fix $K$. Every translation-invariant Gibbs measure $\mu$ of the interaction $\Phi_K$ has decaying covariances on $\mathfrak F$: $\operatorname{Cov}_\mu(w_1,\tau_aw_2)\to0$ as $|a|\to\infty$ for every pair of disjoint bumps, with $w_i=\partial_{t_i}\Delta_\infty(0,\cdot)$. The exponential form asks for $|\operatorname{Cov}_\mu(w_1,\tau_aw_2)|\le C(f_1,f_2)\,e^{-m|a|}$ with some $m>0$.
\end{definition}

\paragraph{Where it enters.} The link is the following theorem.

\begin{theorem}[Theorem W1$_{\mathrm{fin}}$ and Corollary W2$_J$ of \file{SOURCES_FH_U1_REPAIR.md} \S3.2--\S3.3 \status{proved given L3 and R-042; unverified}]\label{thm:open-W1}
For all $K\ge K_0$, all even $\ell>22M$, all disjoint bumps $f_1,f_2$ of diameter $\le D$, and all $a$ with $4D\le|a|\le\ell/2-2D$,
\[
\bigl|\kappa_2^{(K,\ell)}(f_1,\tau_af_2)-\operatorname{Cov}_{\nu_{K,\ell}}(w_1,w_{2,a})\bigr|
\le C_W\,N(f_1)N(f_2)\,(1+|\operatorname{supp}f_1|_1)\,e^{-c_W|a|/M},
\]
with $C_W$, $c_W$ free of $K$, $\ell$, $a$, and the position of the bumps ($C_W$ depends on $D/M$). Hence, for every joint limit $(K_i,\ell_i)$ of Theorem IV with limit two-point cumulant $\kappa_2^J$, and for every bump pair: (i) $\kappa_2^J(f_1,\tau_af_2)\to0$ as $|a|\to\infty$ if and only if $\lim_{|a|\to\infty}\limsup_i|\operatorname{Cov}_{\nu_{K_i,\ell_i}}(w_1,w_{2,a})|=0$; (ii) for $m<c_W/M$, exponential decay at rate $m$ of one is equivalent to exponential decay at rate $m$ of the other.
\end{theorem}

The inputs are L3 and the common radius $c_*r$ of R-042; L6 and L14 are not used. \file{FABLE_HARD.md} \S4.4 states its Theorem W1 ``in every Theorem IV limit'' with inputs L3, L6, L14, but its objects $W_\infty$, $\Delta_\infty$, $\nu_\infty$ are built at fixed $K$. The repair file proves the finite-volume form above instead. Its consequences are these (\file{SOURCES_FH_U1_REPAIR.md} \S3.5, unverified).
\begin{enumerate}
\item Fixed $K$ \status{proved given L3, L14, (Q), R-042; unverified}. If (Gap$_K$) holds, every subsequential infinite-volume lattice limit $\ell_j\to\infty$ at that $K$ has a clustering plaquette two-point function, exponential at rate $m<c_W/M$ under the exponential form. This is a lattice statement at spacing $\varepsilon=L^{-K}$, not a continuum statement. It is the correct content of the sentence of \file{FABLE_HARD.md} \S4.4 ``(Gap$_K$) $\Rightarrow$ OS4 for every Theorem IV limit''.
\item Joint limits \status{proved given L3, R-042; unverified}. If (Gap$_U$) holds, every limit family of Theorem IV satisfies OS4 for the plaquette two-point function. Under the exponential form at rate $m$, the limit has a mass gap $\ge\min(m,c_W/M)$ on the subspace generated by $\{w_f\Omega\}$. The $n$-point form is a sketch.
\item (Gap$_K$) for every $K$ is not shown to suffice for a joint limit, even with decay uniform in $K$. A joint sequence uses each $\Phi_K$ at one finite $\ell$ only, and no file gives a $K$-uniform rate for the approach of $\nu_{K,\ell}$ to the Gibbs states of $\Phi_K$. A second sufficient condition is (K-conv) with (Gap$_\infty$): $\Phi_K\to\Phi_\infty$ in the quasi-local norm with convergent responses, and decay for the Gibbs states of $\Phi_\infty$ \status{sketch, given (K-conv) and (Q) uniformly in $K$}. (K-conv) is an E13-type ultraviolet statement, open.
\item Necessity. By Corollary W2$_J$, OS4 for one joint limit is equivalent to decay of $\operatorname{Cov}_{\nu_{K_i,\ell_i}}(w_1,w_{2,a})$ along its own sequence. Any proof of OS4 for joint limits along the route of \file{INFINITE_VOLUME_FABLE.md} (3.2) must prove this decay for the measures that occur.
\end{enumerate}
Theorem IV (L-196) lists OS4 and the mass gap as not asserted, and Lemma L17 is graded \status{blocked: infrared}. OS1-R (full O(4)) is a separate obstruction, and neither (Gap$_U$) nor (Gap$_K$) addresses it (\file{FABLE_HARD.md} \S2.4).

\paragraph{What is known.} \file{FABLE_HARD.md} \S4.5 calls (Gap$_K$) ``the wall''. (Gap$_K$) and (Gap$_U$) are correlation-decay statements for a weakly coupled lattice gauge theory on $\mathbb Z^4$ at coupling $g\le\gamma$, with no small parameter (for (Gap$_U$), uniformly in the level $K$ and the side $\ell$): the expected decay rate $1/\xi_g$ is exponentially small in $g^{-2}$. The known facts are these.
\begin{enumerate}
\item Why the Step-B obstruction blocks the present method. Step~B of (FX) is a cluster expansion of the unit-scale small-field integral $\int dV_K\,\chi\exp(A^f_K+\sum_XR_{\mathrm{fin}}(X))$ on the whole torus $T_\ell$ at coupling $g$, with no block constraint (\file{INFINITE_VOLUME_FABLE.md} \S2.1). If its activities decayed exponentially with $\ell$-free constants, all connected unit-scale correlations would decay at rate $\kappa/M$, that is, the theory would have a unit-scale mass gap of order one at $g\le\gamma$. Asymptotic freedom predicts a correlation length $\xi_g\sim e^{c/g^2}\gg1$. So Step~B cannot hold with $\ell$-free constants \status{I}. At fixed $\ell$ it holds with constants that grow with $\ell$. An $\ell$-uniform Step~B was one possible way to (Gap$_U$); by this argument it is not available, and Corollary W2$_J$ shows that every proof of OS4 for joint limits on this route must prove the covariance decay of Theorem~\ref{thm:open-W1} along the sequence instead.
\item The finite window. Theorem FD (\S\ref{sec:open-ulap}, item 4) gives Step~B with $\ell$-free constants for $\ell\le22ML\,\xi_g$. Beyond the window the first inequality to fail is the Kotecký--Preiss/Cammarota smallness $\exp(-p_0(g_{K+k}))\le\varepsilon_{\mathrm{KP}}(d,L,M)$ at depth $k>k_*(g)$. It fails as an inequality, not as a divergence of constants: no statement with any constant is proved at $g>\gamma$, and the small-field set is empty as $g\to1$. Continuation of the renormalization group past the unit scale is graded \status{refuted as a route to $\ell\to\infty$; proved as a window statement}.
\item Reflection positivity alone. For the transfer matrix $T=e^{-H}$ of any Theorem IV limit and a bounded observable $A$ of the half-space algebra, $\langle\Theta(A)\tau_tA\rangle=\int_{[0,1]}s^t\,d\mu_A(s)$, so the time-direction two-point function is non-increasing, log-convex, and convergent as $t\to\infty$, with limit $\|P_1A\Omega\|^2$ (Proposition RP1, \status{proved, standard}). Clustering for $A$ is equivalent to $P_1A\Omega=\langle\Omega,A\Omega\rangle\Omega$, that is, to uniqueness of the vacuum, which is the open statement. Chessboard and Gaussian-domination bounds give no decay. Griffiths and FKG inequalities are not available for SU(2). This route is graded \status{refuted as a replacement for the cluster expansion}.
\item Ergodic decomposition. Each ergodic component of $\nu_\infty$ clusters in the mean on $\mathfrak F$. With Proposition RP1 this would give pointwise clustering in the time direction if the component were itself reflection positive. Reflection positivity of an arbitrary ergodic Gibbs component is not known; \file{FABLE_HARD.md} \S4.3 marks this as unsettled.
\item Regime separation. The parameter range of Ba{\l}aban's construction at the unit scale and the strong-coupling range where a lattice mass gap is proved are separated by at least seventeen orders of magnitude in $\beta_W$ (Section~\ref{sec:rem-gap}). The project files contain no numerical test of (Gap$_U$), (Gap$_K$), or (U-Lap)$_K$; \file{FABLE_HARD.md} used no numerics.
\end{enumerate}

\begin{remark}[equivalent, stronger, and weaker forms of (Gap$_U$) and (Gap$_K$)]\label{rem:open-gap-forms}
\begin{enumerate}
\item Equivalent form for one joint limit. For one joint limit sequence, OS4 for the plaquette two-point function is equivalent to decay of $\operatorname{Cov}_{\nu_{K_i,\ell_i}}(w_1,w_{2,a})$ along that sequence, and a mass gap $m<c_W/M$ is equivalent to exponential decay at rate $m$ (Theorem~\ref{thm:open-W1}). This concerns the measures that occur along the sequence, not all Gibbs states of any $\Phi_K$. (Gap$_U$) asks for this decay along every sequence at once.
\item (Gap$_U$) and (Gap$_K$). (Gap$_U$) concerns the finite-volume measures $\nu_{K,\ell}$, uniformly in $(K,\ell)$. (Gap$_K$) concerns the infinite-volume Gibbs states of $\Phi_K$ at one $K$. The sources prove no implication between them. (Gap$_K$) for every $K$ gives clustering at each fixed $K$ only (item 1 after Theorem~\ref{thm:open-W1}).
\item Why (Gap$_K$) quantifies over all states. Different Gibbs states of $\Phi_K$ can have different covariance decay. An argument that uses only $(W_\infty,\Delta_\infty)$ and bounds valid for every state cannot prove clustering at fixed $K$ (\file{FABLE_HARD.md} \S4.4). A proof must either identify the specific limit state or prove decay for all translation-invariant Gibbs states; (Gap$_K$) is the second option.
\item A possibly weaker form. Both conditions ask for decay of first-order responses only. Whether this is strictly easier than decay for all local functionals is not settled; $\mathfrak F$ is the family of sourced response functionals and need not separate Gibbs states (Remark~\ref{rem:open-ulap-forms}(3)).
\item A stronger form. \file{PLAN.md} (entry for \file{ALT_ROUTES.md}) restates the common target of three routes as boundary-condition-uniform and volume-uniform exponential decay of the unit-scale measure at distance $\xi_g$, and calls it ``= (Gap$_K$)''. That statement is stronger than (Gap$_K$) as defined in \file{FABLE_HARD.md} \S4.4, which asks for decay on $\mathfrak F$ of translation-invariant Gibbs states only, with no rate in its plain form. This draft does not compare it with (Gap$_U$). \file{ALT_ROUTES.md} has no verifier report.
\item Relation to (U-Lap)$_K$. The sources give no implication between (Gap$_K$) or (Gap$_U$) and (U-Lap)$_K$. All three are infrared statements about the family $\nu_{K,\ell}$, and \file{FABLE_HARD.md} \S4.3 links clustering to vacuum uniqueness: Proposition RP1 reduces clustering to vacuum uniqueness, which the source calls ``(U-Lap)/(Gap) again''.
\end{enumerate}
\end{remark}

\subsection{Summary of the two problems}\label{sec:open-summary}

\begin{center}
\begin{tabular}{p{0.15\textwidth}p{0.37\textwidth}p{0.38\textwidth}}
\hline
problem & if it holds & status in the sources \\
\hline
(U-Lap)$_K$ & at fixed $K$, the $\ell\to\infty$ limit of Theorem IV exists along the full sequence (Theorem U1$'$(c)); the joint limit needs E13-type UV uniqueness and uniformity in $K$ in addition & open; equivalent to (U-IV$^=$)$_K$ and sufficient for (U-IV)$_K$, given (L6-$\mathcal P$), L3, L14 \status{proved, unverified}; implied by (U-Gibbs)$_K$; no proof route \\
(Gap$_U$) & OS4 for the plaquette two-point function of every joint limit of Theorem IV; under the exponential form at rate $m$, a mass gap $\ge\min(m,c_W/M)$ on the span of $\{w_f\Omega\}$ & open; given L3 and R-042 \status{proved, unverified} (Theorem W1$_{\mathrm{fin}}$, Corollary W2$_J$); (Gap$_K$) for every $K$ gives clustering only at each fixed $K$ and is not shown to suffice for joint limits; $\ell$-uniform Step~B, RP alone, and RG continuation are \status{refuted} as routes \\
\hline
\end{tabular}
\end{center}
